\noindent This appendix is designed as a last-minute revision sheet. The formulas matter, but the real skill is knowing when each one applies.

\subsubsection*{Event Probability}
If all outcomes are equally likely, then
\[
P(A)=\frac{n(A)}{n(S)}.
\]

\subsubsection*{Complement Rule}
\[
P(A^c)=1-P(A).
\]

\subsubsection*{Addition Rule}
\[
P(A\cup B)=P(A)+P(B)-P(A\cap B).
\]
If $A$ and $B$ are mutually exclusive, then
\[
P(A\cup B)=P(A)+P(B).
\]

\subsubsection*{Conditional Probability}
\[
P(A\mid B)=\frac{P(A\cap B)}{P(B)},
\qquad P(B)>0.
\]

\subsubsection*{Multiplication Rule}
\[
P(A\cap B)=P(A)P(B\mid A)=P(B)P(A\mid B).
\]

\subsubsection*{Independence}
If $A$ and $B$ are independent, then
\[
P(A\cap B)=P(A)P(B).
\]

\subsubsection*{Expectation and Variance}
For a discrete random variable $X$,
\[
E(X)=\sum xP(X=x),
\]
\[
E(X^2)=\sum x^2P(X=x),
\]
\[
\mathrm{Var}(X)=E(X^2)-[E(X)]^2.
\]

\subsubsection*{Binomial Distribution}
If
\[
X\sim \mathrm{Bin}(n,p),
\]
then
\[
P(X=r)=\binom{n}{r}p^r(1-p)^{n-r},
\qquad r=0,1,\dots,n,
\]
\[
E(X)=np,
\qquad
\mathrm{Var}(X)=np(1-p).
\]

\subsubsection*{Counting Formulas}
\[
{}^nP_r=\frac{n!}{(n-r)!},
\qquad
\binom{n}{r}=\frac{n!}{r!(n-r)!}.
\]

\subsubsection*{When to Use What}
\begin{center}
\begin{tabular}{p{5cm} p{9cm}}
\toprule
\textbf{Question wording} & \textbf{Likely method} \\
\midrule
at least one & complement rule \\
given that & conditional probability \\
without replacement & multiplication rule with changing fractions \\
committee or selection & combinations \\
exactly $r$ successes in repeated trials & binomial formula \\
fair game or average return & expected value \\
\bottomrule
\end{tabular}
\end{center}
